\documentclass{article}
% Packages
\usepackage[utf8]{inputenc}

\usepackage{amsmath}
\usepackage{amssymb}
\usepackage{tabularx}
\usepackage{xcolor}
\usepackage{eurosym}

\usepackage{geometry}
\geometry{a4paper, total={160mm,240mm}, left=25mm, top=20mm}

%\usepackage[citestyle=authoryear, maxcitenames = 2, hyperref=true]{biblatex} %Imports biblatex package
\usepackage[citestyle=authoryear, maxcitenames = 2, hyperref=true, backref=false, abbreviate=false, maxbibnames=10]{biblatex}
\addbibresource{references.bib}
\usepackage{hyperref}
\usepackage{soul}

\newcommand{\R}{\mathbb{R}}
\newcommand{\N}{\mathbb{N}}
\newcommand{\E}{\mathbb{E}}

\newcommand{\SOC}{\text{SoC}}
\newcommand{\soc}{\text{soc}}

\newtheorem{notation}{Notation}[section]
\newtheorem{note}{Note}[section]

\setlength{\parindent}{0pt}

\title{EC Model Formulation - Intraday Market 2}
\author{F.- Javier Heredia Cervera, Ignasi Mañé, Albert Sol\`a Vilalta, Jan Jettmann}
\date{\today}



\begin{document}

\maketitle

\section{Introduction}
This model is based on the implementation of the DA model. This model allows the EC to place bids into the second intraday market, using information that was revealed after the fist IM has cleared. The results of the DA, RM and IM1 are taken as parameters. This model can be included in the modular architecture that allows the simulation of the EC participating in several markets by placing bids, using the most up-to-date information at that time.\\
We consider an energy community with the following elements:
\begin{enumerate}
	\item (Implicit) Flexible Demand.
	\item Wind Farm.
	\item Solar Farm.
	\item Battery Energy Storage System (BESS).
\end{enumerate}

In the following sections we present the proposed problem formulation. Since it is quite a large model, the sets, parameters, variables, and constraints are defined on separate sections.

\begin{notation}
	In the definition of sets, parameters, and variables, we use parenthesis to denote the name of that element in the code. For example, we use $\mathcal{T}$ (\verb|T|) to denote that the set $\mathcal{T}$ has the name \verb|T| in the code.  
\end{notation}

\begin{note} \label{NOTE:ScenarioTreatment}
	We do not distinguish in this formulation between ``complete set of scenarios'' and ``preserved set of scenarios''. This is a pre-processing step, and does not affect the formulation. 
\end{note}

\section{Fundamental Elements}
\label{SEC:FundamentalElements}

\begin{tabularx}{\textwidth}{l l}
    \textbf{Sets:} \\
    $\mathcal{T}$ (\verb|T|) & set of $T$ time periods $\{ 1, \dots, T \}$ [hours]. \\
	$\mathcal{T}_0$ (\verb|T0|) & extended set of time periods $\{ 0, \dots, T \}$ [hours]. \\     
	$\mathcal{I}$ (\verb|IM|) & set of $I$ intraday markets $\{1, \dots, I \}$. \\    
    $\mathcal{T}_i$ (\verb|TIM|) & set of time periods in intraday market $i \in \mathcal{I}$ [hours]. \\
	$\mathcal{I}_t$ (\verb|IMT|) & set of intraday markets where $t \in \mathcal{T}$ is a time period. \\    
	$\Omega$ (\verb|S|) & set of scenarios. \\
	$\mathcal{SG}$ (\verb|SG|) & set of stages $\{1, \dots, SG \}$. \\
	$\mathcal{SG}_0$ (\verb|SG0|) & extended set of stages $\{0, \dots, SG \}$. \\     
    $\mathcal{C}_{sg, \omega}$ (\verb|c|) & scenario cluster of scenario $\omega \in \Omega$ at stage $sg \in \mathcal{SG}_0$. \\
    $\mathcal{RV}$ & set of random variables $\{1, \dots, nRV \}$. \\
    \\
   
    
    \textbf{Parameters:} \\
	$T$ (\verb|nT|) & number of time periods. \\
	$I$ (\verb|nIM|) & number of intraday markets. \\
	$nSG$ (\verb|nSG|) & number of stages. \\
	$SG^{IM}_i$ (\verb|sgim|) & stage in which the price of intraday market $i \in \mathcal{I}$ is revealed. \\
	$nRV$ (\verb|nRV|) & number of random variables. \\
	$nRVSG_{sg}$ (\verb|nRVSG|) & number of random variables revealed at stage $sg \in \mathcal{SG}$. \\
	$\Pi_{\omega}$ (\verb|Prob|) & probability of scenario $\omega \in \Omega$. \\
	$S_{rv, \omega}$ (\verb|Scen|) & value of random variable $rv \in \mathcal{RV}$ in scenario $\omega \in \Omega$. \\    
    \\
       
    \textbf{Variables:} \\
	\\
\end{tabularx}

\section{Flexible Demand}
\label{SEC:FlexibleDemand}

\begin{tabularx}{\textwidth}{l l}
    \textbf{Sets:} \\
    $\mathcal{F}$ (\verb|FI|) & set of time intervals where a fraction of demand $c_f$ has to be met. \\ 
    \\
    
    \textbf{Parameters:} \\
    	$nF$ (\verb|nFI|) & number of intervals in $\mathcal{F}$. \\
    $D_t$ (\verb|FD|) & central demand at $t \in \mathcal{T}$ [MWh]. \\
    $\underline{D}_t$ (\verb|FD_L|) & minimum demand at $t \in \mathcal{T}$ (inflexible demand) [MWh]. \\
    $\overline{D}_t$ (\verb|FD_U|) & maximum demand at $t \in \mathcal{T}$ [MWh]. \\
    $T_f^-$ (\verb|TF_L|) & first time step of interval $f \in \mathcal{F}$ [hour]. \\
    $T_f^+$ (\verb|TF_U|) & last time step of interval $f \in \mathcal{F}$ [hour]. \\
    $c_f$ (\verb|coef_FD|) & fraction of demand that needs to be consumed during interval $f \in \mathcal{F}$. \\
    $C^{FD}$ (\verb|C_FD|) & cost per unit of flexible demand used [\euro/MWh]. \\
    \\
       
    \textbf{Variables:} \\
    $fd_{t, \omega}$ (\verb|var_fd|) & flexible demand at $t \in \mathcal{T}$ under scenario $\omega \in \Omega$ (after flexibility is \\
    & considered) [MWh]. \\
    $afd_{t, \omega}^+$ (\verb|var_afd_p|) & auxiliary variable. Positive part of flexible demand shift from central \\
    & demand $D_t$ [MWh]. \\
    $afd_{t, \omega}^-$ (\verb|var_afd_m|) & auxiliary variable. Negative part of flexible demand shift from central \\
    & demand $D_t$ [MWh]. \\
    \\
    \end{tabularx}

First, we introduce nonnegativity constraints to the flexible demand variables (\verb|at variable| \verb|definition|)
\begin{equation} \label{EQ:FlexDemandPos}
	fd_{t, \omega} \geq 0, \ afd_{t, \omega}^+ \geq 0, \ afd_{t, \omega}^- \geq 0 \quad \forall t \in \mathcal{T}, \ \forall \omega \in \Omega.
\end{equation}

The flexible demand is bounded by the minimum demand and maximum demand at every time step. This is
\begin{equation} \label{EQ:FlexDemand_UB_preliminary}
	fd_{t, \omega} \leq \overline{D}_t \quad \forall t \in \mathcal{T}, \ \forall \omega \in \Omega
\end{equation}
and 
\begin{equation} \label{EQ:FlexDemand_LB_preliminary}
	\underline{D}_t \leq fd_{t, \omega} \quad \forall t \in \mathcal{T}, \ \forall \omega \in \Omega.
\end{equation}

Moreover, the total daily demand has to meet the central estimate of the demand over the day (\verb|DailyDem|): 
\begin{equation} \label{EQ:DailyDem}
	\sum_{t \in \mathcal{T}} fd_{t, \omega} = \sum_{t \in \mathcal{T}} D_t \quad \forall \omega \in \Omega.
\end{equation}
This is done to avoid meeting a lower demand over the full time horizon of the problem ($24$h), or a higher demand if the reward for downward reserve is consistently high and/or electricity prices are consistently negative. We also introduce time intervals where a certain percentage of the central demand needs to be met over that interval (\verb|InterDem|):
\begin{equation} \label{EQ:InterDem}
	\sum_{T_f^- \leq t \leq T_f^+} fd_{t, \omega} \geq c_f \sum_{T_f^- \leq t \leq T_f^+} D_t \quad \forall f \in \mathcal{F}, \ \forall \omega \in \Omega.
\end{equation}

To model the cost of implicit flexible demand, we introduce a penalty in the objective function to use it. This requires to split the variables $fd_{t, \omega}$ into its positive and negative parts, which is done in the following constraints (\verb|FlexDemDisplace|)
\begin{equation} \label{EQ:FlexDemDisplace}
	D_t - fd_{t, \omega} = afd_{t, \omega}^+ - afd_{t, \omega}^- \quad \forall t \in \mathcal{T}, \ \forall \omega \in \Omega.
\end{equation}



\section{Wind Farm}
\label{SEC:WindFarm}

\begin{tabularx}{\textwidth}{l l}
    \textbf{Sets:} \\
    \\
    
    \textbf{Parameters:} \\
	$SG^{WP}_t$ (\verb|sgpw|) & stage in which the wind production at time $t \in \mathcal{T}$ is revealed. \\
	$W_{t, \omega}$ (\verb|pW|) & wind production at time $t \in \mathcal{T}$ under scenario $\omega \in \Omega$ [MWh]. \\
	$\overline{P}^W$ (\verb|max_pW|) & wind nameplate capacity [MW]. \\
	\\  
\end{tabularx}

\section{Solar Farm}
\label{SEC:SolarFarm}

\begin{tabularx}{\textwidth}{l l}
    \textbf{Sets:} \\
    \\
    
    \textbf{Parameters:} \\
	$PV_{t, \omega}$ (\verb|pPV|) & solar production at time $t \in \mathcal{T}$ under scenario $\omega \in \Omega$ [MWh]. \\
	$\overline{P}^{PV}$ (\verb|max_pPV|) & solar nameplate capacity [MW]. \\
	\\  
\end{tabularx}

We assume that observations of the solar random variable at time $t \in \mathcal{T}$ is made at the same stage as the wind random variable at time $t \in \mathcal{T}$.

\section{Battery Energy Storage System}
\label{SEC:BatteryEnergyStorageSystem}

\begin{tabularx}{\textwidth}{l l}
    \textbf{Sets:} \\
 	
    \\
    
    \textbf{Parameters:} \\
	$E^B$ (\verb|Emax|) & BESS energy capacity [MWh].	 \\
	$P^B$ (\verb|Dmax|) & BESS maximum charging/discharging rate [MW]. \\
	$\eta^B$ (\verb|RTE|) & BESS round-trip efficiency. \\
	$\overline{\SOC}$ (\verb|SOCmax|) & BESS maximum state of charge. Unitless, values in $[0, 1]$. \\
	$\underline{\SOC}$ (\verb|SOCmin|) & BESS minimum state of charge. Unitless, values in $[0, 1]$. \\
	$\SOC_0$ (\verb|SOCini|) & BESS initial state of charge. Unitless, values in $[0, 1]$. \\
	$\SOC_T$ (\verb|SOCfin|) & BESS final state of charge. Unitless, values in $[0, 1]$. \\
	\\
       
    \textbf{Variables:} \\
	$c_{t, \omega}$ (\verb|cV|)& BESS charge rate at time $t \in \mathcal{T}$ under scenario $\omega \in \Omega$ [MW]. \\
	$d_{t, \omega}$ (\verb|dV|)& BESS discharge rate at time $t \in \mathcal{T}$ under scenario $\omega \in \Omega$ [MW]. \\	
	$id_{t, \omega}$ (\verb|idV|) & binary. Indicates the charging ($id_{t, \omega} = 0$) or discharging ($id_{t, \omega}= 1$) state \\
	& of the BESS at time $t \in \mathcal{T}$ under scenario $\omega \in \Omega$. \\
	$\soc_{t, \omega}$ (\verb|socV|) & state of charge of BESS at time $t \in \mathcal{T}_0$ under scenario $\omega \in \Omega$. Unitless, \\
	& values in $[0, 1]$. \\
	\\
\end{tabularx}

First, we introduce the following nonnegativity  constraints, and define $id_{t, \omega}$ as a binary (\verb|at variable| \verb|definition|):
\begin{equation} \label{dVdCPos}
	c_{t, \omega} \geq 0, \quad d_{t, \omega} \geq 0, \quad id_{t, \omega} \in \{ 0, 1 \} \quad \forall t \in \mathcal{T}, \ \forall \omega \in \Omega.
\end{equation}

We introduce limits in the charged and discharged quantities of the battery, as well as not allowing simultaneous charging and discharging in every time step (\verb|VPP_state_dV|):
\begin{equation} \label{EQ:VPP_state_dV}
	d_{t, \omega} \leq P^B id_{t, \omega} \quad \forall t \in \mathcal{T}, \ \forall \omega \in \Omega
\end{equation}
and (\verb|VPP_state_cV|)
\begin{equation} \label{EQ:VPP_state_cV}
	c_{t, \omega} \leq P^B (1 - id_{t, \omega}) \quad \forall t \in \mathcal{T}, \ \forall \omega \in \Omega.
\end{equation}

We define the variable state of charge $\soc_{t, s}$ of the BESS with constraints (\verb|SOCV|): 
\begin{equation} \label{EQ:SOCV}
	\soc_{t, \omega} = \soc_{t - 1, \omega} + \frac{c_{t, \omega} - \frac{d_{t, \omega}}{\eta^B}}{E^B} \quad \forall t \in \mathcal{T}, \ \forall \omega \in \Omega. 
\end{equation} 
They have to be within the operational limits of the BESS (\verb|at variable definition|):
\begin{equation} \label{EQ:socVOperLimits}
	\underline{\SOC} \leq \soc_{t, \omega} \leq \overline{\SOC} \quad \forall t \in \mathcal{T}, \ \forall \omega \in \Omega.
\end{equation}
The initial and final states of charge of the BESS are defined with (\verb|SOCV_ini|):
\begin{equation} \label{EQ:SOCV_ini}
	\soc_{0, \omega} = \SOC_0 \quad \forall \omega \in \Omega
\end{equation}
and (\verb|SOCV_fin|)
\begin{equation} \label{EQ:SOCV_fin}
	\soc_{T, \omega} = \SOC_T \quad \forall \omega \in \Omega.
\end{equation}



\section{Market Participation}
\label{SEC:MarketParticipation}

\subsection{Day Ahead Market}




\begin{tabularx}{\textwidth}{l l}

	\\
    \textbf{Parameters:} \\
	$\lambda^{D}_{t}$ (\verb|lD|) & observed  day-ahead market price at time $t \in \mathcal{T}$  [\euro/MWh]. \\
       \\
	$e^{DA+}_{t}$ (\verb|eDA_p|) & energy matched in the day-ahead market of a selling bid by EC at time $t \in \mathcal{T}$ \\
    \\
	$ $ [MWh]. \\
    
	$e^{DA-}_{t}$ (\verb|eDA_m|) & energy matched in the day-ahead market of a buying bid by EC at time $t \in \mathcal{T}$ \\
	& [MWh]. 
    \\
    $ie^{DA+}_{t,}$ (\verb|ieDA_p|) & binary. $ie^{DA+}_{t} = 1$ if the EC sells electricity at time $t \in \mathcal{T}$ \\
	& $ie^{DA+}_{t} = 0$ o/w. \\
	$ie^{DA-}_{t}$ (\verb|ieDA_m|) & binary. $ie^{DA-}_{t} = 1$ if the EC buys electricity at time $t \in \mathcal{T}$ \\
	& $ie^{DA-}_{t} = 0$ o/w.
\end{tabularx}

Since the values of these parameters are obtained from running the DA model, it is assumed that they obey to all constraints formulated in that model, thus they do not need to be listed in the IM model as well.

\subsection{Reserve Market}
We assume that the wind and solar farms cannot contribute to provide reserve. The other elements of the energy community (Battery and Flexible Demand) can. This follows the approach by \textcite{Heredia18}, where VPP's bid to reserve markets is limited by the battery capabilities only.

%To model reserve bids, we use recourse variables only because price-accepting reserve bids do not seem to have a clear motivation. From these recourse variables, we build buying and selling bid curves a posteriori, from the results obtained at each scenario. The role of first stage variables is represented in the reserve by the binary variable indicating if the EC submitted a bid to the day-ahead market. Recall that  MIBEL market rules establish that a market participant needs to have an accepted day-ahead market bid to participate in the reserve market.

\begin{tabularx}{\textwidth}{l l}

\textbf{Parameters:} \\
$\lambda^{R}_{t,}$ (\verb|lR|) & observed reserve market price at time $t \in \mathcal{T}$ [\euro/MWh]. \\
$T^R$ (\verb|TSR|) & duration of reserve response [h]. \\
$\lambda^{R,\text{undel}}_{t,\omega}$ (\verb|lR_penalty|) & penalty price applied to undelivered reserve at time $t \in \mathcal{T}$ under scenario $\omega \in \Omega$ [\euro/MWh]. \\    
$r^{U}_{t, \omega}$ (\verb|rU|) & upward reserve matched by the EC at time $t \in \mathcal{T}$ under scenario $\omega \in \Omega$ [MW]. \\
$r^{D}_{t, \omega}$ (\verb|rD|) & downward reserve matched by the EC at time $t \in \mathcal{T}$ under scenario $\omega \in \Omega$ [MW]. \\
\\
\textbf{Variables:} \\
$r^{U, B}_{t, \omega}$ (\verb|rU_B|) & BESS component of upward reserve matched at time $t \in \mathcal{T}$ \\
& under scenario $\omega \in \Omega$ [MW]. \\
$r^{D, B}_{t, \omega}$ (\verb|rD_B|) & BESS component of downward reserve matched at time $t \in \mathcal{T}$ \\
& under scenario $\omega \in \Omega$ [MW]. \\
$r^{U, FD}_{t, \omega}$ (\verb|rU_FD|) & flexible demand component of upward reserve matched at time $t \in \mathcal{T}$ \\
& under scenario $\omega \in \Omega$ [MW]. \\
$r^{D, FD}_{t, \omega}$ (\verb|rD_FD|) & flexible demand component of downward reserve matched at time $t \in \mathcal{T}$ \\
& under scenario $\omega \in \Omega$ [MW]. \\
\\
$r^{U,\text{undel}}_{t,\omega}$ (\verb|rU_undelivered|) & portion of the upward reserve contracted at time $t\in\mathcal{T}$ that cannot be delivered\\ in scenario $\omega\in\Omega$ [MW]. \\
\\
$r^{D,\text{undel}}_{t,\omega}$ (\verb|rD_undelivered|) & portion of the downward reserve contracted at time $t\in\mathcal{T}$ that cannot be delivered\\ in scenario $\omega\in\Omega$ [MW]. \\

\end{tabularx}\\
\\
\\
While the total amount of reserve the EC bid is fixed, the percentages of flexible demand and BESS that will provide it are left to optimization. The idea is that by leaving the BESS and FD parts of the reserve as variables, the EC has greater flexibility on how it will honor the reserve bid depending on the scenario. The variables are kept open up to the stage that they are revealed.\\
First, we impose that the reserve variables are nonnegative (\verb|at variable definition|)
\begin{equation} \label{EQ:RVarDef}
    r^{U}_{t, \omega} \geq 0, \quad r^{D}_{t, \omega} \geq 0, \quad r^{U, B}_{t, \omega} \geq 0, \quad r^{D, B}_{t, \omega} \geq 0, \quad r^{U, FD}_{t, \omega} \geq 0, \quad r^{D, FD}_{t, \omega} \geq 0, r^{U,\text{undel}}_{t,\omega} \geq 0, \quad r^{D,\text{undel}}_{t,\omega} \geq 0, 
\end{equation}

for all $t \in \mathcal{T}$, and $\omega \in \Omega$. 


The reserve market parameters mentioned before already  abide to the constraints imposed in the previous stages models, so we impose that the penalty reserve variables are nonnegative (\verb|at variable definition|) for the new variables.\\
\begin{equation} \label{EQ:RVarDef}
    r^{U,\text{undel}}_{t,\omega} \geq 0, \quad r^{D,\text{undel}}_{t,\omega} \geq 0, 
\end{equation}
for all $t \in \mathcal{T}$, and $\omega \in \Omega$. 

The reserve provided by the EC is the sum of the reserve provided by those elements able to provide reserve and the undelivered reserve. For various reasons it can happen that using the results of the reserve market model run as parameters for the IM model, infeasibilities happen for which the EC cannot deliver the contracted reserve.  \\
This is (\verb|RM_components_U|)\\
\begin{equation} \label{EQ:RUcom}
    r^{U}_{t,\omega} \;=\;
        r^{U,B}_{t,\omega}
      \;+\; r^{U,FD}_{t,\omega}
      \;+\; r^{U,\text{undel}}_{t,\omega},
    \qquad \forall\, t\in\mathcal T,\;\omega\in\Omega ,
\end{equation}

and for downward reserve (\verb|RM_components_D|)

\begin{equation} \label{EQ:RDcomp}
    r^{D}_{t,\omega} \;=\;
        r^{D,B}_{t,\omega}
      \;+\; r^{D,FD}_{t,\omega}
      \;+\; r^{D,\text{undel}}_{t,\omega},
    \qquad \forall\, t\in\mathcal T,\;\omega\in\Omega .
\end{equation}

\subsection{Intraday Markets}

The intraday market model follows the same structure as the previous ones, with an ordered set that allows for the creation of bid curves.

\begin{tabularx}{\textwidth}{l l}
    \textbf{Sets:} \\
    	$\mathcal{SSI}_2$ (\verb|Ssi_2|) & $\{(t, l, j) \in \mathcal{T}_2 \times \Omega \times \Omega \hspace{2mm} | \hspace{2mm} l \neq j, \hspace{2mm} \lambda^I_{2, t, l} \leq \lambda^I_{2, t, j}, \hspace{2mm}$ \\ 
	& Rather technical. Used to ensure monotonicity of intraday market bids \\
	& w.r.t. intraday market price realisations. \\
	& For $t \in \mathcal{T}_2$ provides the scenario pairs $(l, j)$ \\
	& s.t. $\lambda^I_{2, t, l} \leq \lambda^I_{2, t, j}$ (order is important).\\
    & Notice that this set is only for the second intraday market, since the other ones are,\\
    & from this model's view, decision variables.
    \\
    
    \textbf{Parameters:} \\
    	$\lambda^I_{i, t, \omega}$ (\verb|lI|) & intraday market $i \in \mathcal{I}$ price at time $t \in \mathcal{T}$ under scenario $\omega \in \Omega$ [\euro/MWh]. \\
            $e^{IM}_{1, t, \omega}$ (\verb|eIM|) & energy matched by the EC at intraday market $1 \in \mathcal{I}$ at time $t \in \mathcal{T}$ \\
    & under scenario $\omega \in \Omega$ [MWh]. \\

	$R^{IDA}$ (\verb|maxTIM|) & bound on the ratio of energy traded in intraday markets and day ahead \\
	& market (to avoid speculation).\\
    \\ 
       
    \textbf{Variables:} \\
    $e^{IM}_{i, t, \omega}$ (\verb|eIM|) & energy matched by the EC at intraday market $i \in \mathcal I \setminus \{1\}$ at time $t \in \mathcal{T}$ \\
    & under scenario $\omega \in \Omega$ [MWh]. \\
	\\
\end{tabularx}

To avoid speculative behaviour, the bid of the EC to intraday markets should be small compared to the day ahead market (\verb|IM_bounds_1| and \verb|IM_bounds_2|)
\begin{equation} \label{EQ:DM_IM_1}
	-R^{IDA} (e^{DA+}_{t} + e^{DA-}_{t})  \leq \sum_{i \in \mathcal{I}_t} e^{IM}_{i, t, \omega} \leq R^{IDA} (e^{DA+}_{t} + e^{DA-}_{t}) \quad \forall t \in \mathcal{T}, \ \forall \omega \in \Omega
\end{equation}
and (\verb|IM_bounds_3| and \verb|IM_bounds_4|)
\begin{equation} \label{EQ:IM_bounds_3-4}
	-R^{IDA} (e^{DA+}_{t} + e^{DA-}_{t}) \leq e^{IM}_{i, t, \omega} \leq R^{IDA} (e^{DA+}_{t} + e^{DA-}_{t}) \quad \forall i \in \mathcal{I}, \ \forall t \in \mathcal{T}_i, \ \forall \omega \in \Omega.
\end{equation}

To build a bidding curve for the second intraday market from the results obtained at each scenario, the quantities matched in intraday markets have to be monotonically increasing w.r.t. the scenario prices. This is guaranteed with the following constraints (\verb|IM_bid_mono_2_rule|):
\begin{equation} \label{EQ:lI_bid}
	e^{IM}_{2, t, \omega} \leq e^{IM}_{2, t, \omega'} \quad \forall (t, \omega, \omega') \in \mathcal{SSI}.
\end{equation}

\subsection{System Imbalances}



\begin{tabularx}{\textwidth}{l l}
    \textbf{Parameters:} \\
    $\lambda^{IB+}_{t, \omega}$ (\verb|lPIB|) & positive imbalance price at time $t \in \mathcal{T}$ under scenario $\omega \in \Omega$ [\euro/MWh]. \\
	$\lambda^{IB-}_{t, \omega}$ (\verb|lNIB|) & negative imbalance price at time $t \in \mathcal{T}$ under scenario $\omega \in \Omega$ [\euro/MWh]. \\
    $P^{IB+}_{t, \omega}$ (\verb|PIB_p|) & upper bound on the positive imbalance at time $t \in \mathcal{T}$ \\
    & under scenario $\omega \in \Omega$ [MWh]. \\
	$P^{IB-}_{t, \omega}$ (\verb|PIB_m|) & upper bound on the negative imbalance at time $t \in \mathcal{T}$ \\ 
	& under scenario $\omega \in \Omega$ [MWh]. \\
    \\ 
       
    \textbf{Variables:} \\
    $p_{t, \omega}^{IB-}$ (\verb|pIB_m|) & negative imbalance of EC at time $t \in \mathcal{T}$ under scenario $\omega \in \Omega$ [MWh].	\\
	$p_{t, \omega}^{IB+}$ (\verb|pIB_p|) & positive imbalance of EC at time $t \in \mathcal{T}$ under scenario $\omega \in \Omega$ [MWh].	\\ 
    \\
\end{tabularx}

First, we impose that the imbalance variables are nonnegative (\verb|at variable definition|)
\begin{equation} \label{EQ:IBVarDef}
    p_{t, \omega}^{IB+} \geq 0, \quad p_{t, \omega}^{IB-} \geq 0 \quad \forall t \in \mathcal{T}, \ \forall \omega \in \Omega.
\end{equation}

The imbalances of the EC are defined in the following equation (\verb|Imbalances|)
\begin{equation} \label{EQ:IBDef}
    p^{IB+}_{t, \omega} - p^{IB-}_{t, \omega} = e^{DA-}_{t} + W_{t, \omega} + PV_{t, \omega} + d_{t, \omega} - \Big (e^{DA+}_{t} + \sum_{i \in \mathcal{I}_t} e^{IM}_{i, t, \omega} + fd_{t, \omega} + c_{t, \omega} \Big ) \quad \forall t \in \mathcal{T}, \ \forall \omega \in \Omega.
\end{equation}


We introduce upper and lower bounds to the imbalances, to avoid speculative behavior.
Note: This leads to infeasibility issues so the bounds are set to 500 MWh.\\

(\verb|IB_pos_UB|):\\
\begin{equation} \label{EQ:VPP_IB_MAX}
	p_{t, \omega}^{IB+} \leq P_{t, \omega}^{IB+} \quad \forall t \in \mathcal{T}, \ \forall \omega \in \Omega
\end{equation}
and (\verb|IB_neg_UB|)
\begin{equation} \label{EQ:VPP_IB_MIN}
	p_{t, \omega}^{IB-} \leq P_{t, \omega}^{IB-} \quad \forall t \in \mathcal{T}, \ \forall \omega \in \Omega.
\end{equation}

%\begin{equation} \label{EQ:VPP_IB_MAX}
%	p_{t, \omega}^{IB+} \leq P_{t, \omega}^{IB+} \quad \forall t \in \mathcal{T}, \ \forall \omega \in \Omega
%\end{equation}
%and (\verb|IB_neg_UB|)
%\begin{equation} \label{EQ:VPP_IB_MIN}
%	p_{t, \omega}^{IB-} \leq P_{t, \omega}^{IB-} \quad \forall t \in \mathcal{T}, \ \forall \omega \in \Omega.
%\end{equation}



\section{Nonanticipativity Constraints}
\label{SEC:NonanticipativityConstraints}

The nonanticipativity constraints ensure that every decision is made with the amount of information available at the time of making that decision.

For intraday market variables \( e^{IM}_{i,t,\omega} \), non-anticipativity constraints are imposed only for markets where the price is not yet known at the time of bidding, namely \( i = 3 \). For $i = 3$, they are recourse variables. The constraints are defined as follows:

For IM: (\verb|NAC_eIM_2|)
  \begin{equation} \label{EQ:NAC_eIM2}
  	e^{IM}_{2, t, \omega} = e^{IM}_{2, t, \omega'} \quad 
  	\forall t \in \mathcal{T}_2, \ 
  	\forall \omega \in \Omega, \ 
  	\forall \omega' \in \mathcal{C}_{SG^{IM}_2, \omega}, \ \omega \neq \omega'.
  \end{equation}
  and (\verb|NAC_eIM|)
\begin{equation} \label{EQ:NAC_eIM}
	e^{IM}_{3, t, \omega} = e^{IM}_{3, t, \omega'} \quad 
	 \ \forall t \in \mathcal{T}_3, \ 
	\forall \omega \in \Omega, \ 
	\forall \omega' \in \mathcal{C}_{SG^{IM}_3 - 1, \omega}, \ \omega \neq \omega'.
\end{equation}


For imbalance variables, we have (\verb|NAC_pIB_p|):
\begin{equation} \label{EQ:NAC_pPIB}
	p^{IB+}_{t, \omega} = p^{IB+}_{t, \omega'} \quad \forall t \in \mathcal{T}, \ \forall \omega \in \Omega, \ \forall \omega' \in \mathcal{C}_{SG^{WP}_t, \omega}, \ \omega \neq \omega',
\end{equation}
and for negative imbalances (\verb|NAC_pIB_m|):
\begin{equation} \label{EQ:NAC_pNIB}
	p^{IB-}_{t, \omega} = p^{IB-}_{t, \omega'} \quad \forall t \in \mathcal{T}, \ \forall \omega \in \Omega, \ \forall \omega' \in \mathcal{C}_{SG^{WP}_t, \omega}, \ \omega \neq \omega'.
\end{equation}


For flexible demand variables, we have (\verb|NAC_var_fd|)
\begin{equation} \label{EQ:NAC_var_fd}
	fd_{t, \omega} = fd_{t, \omega'} \quad \forall t \in \mathcal{T}, \ \forall \omega \in \Omega, \ \forall \omega' \in \mathcal{C}_{SG^{WP}_t - 1, \omega}, \ \omega \neq \omega',
\end{equation}
and (\verb|NAC_var_afd_p|)
\begin{equation} \label{EQ:NAC_var_afd_p}
	afd^+_{t, \omega} = afd^+_{t, \omega'} \quad \forall t \in \mathcal{T}, \ \forall \omega \in \Omega, \ \forall \omega' \in \mathcal{C}_{SG^{WP}_t - 1, \omega}, \ \omega \neq \omega',
\end{equation}
and (\verb|NAC_var_afd_m|)
\begin{equation} \label{EQ:NAC_var_afd_m}
	afd^-_{t, \omega} = afd^-_{t, \omega'} \quad \forall t \in \mathcal{T}, \ \forall \omega \in \Omega, \ \forall \omega' \in \mathcal{C}_{SG^{WP}_t - 1, \omega}, \ \omega \neq \omega'.
\end{equation}
and (\verb|NAC_rU_FD|)
\begin{equation} \label{EQ:NAC_rUFD}
	r^{U,FD}_{t, \omega} = r^{U,FD}_{t, \omega'} \quad \forall t \in \mathcal{T},\ \forall \omega \in \Omega,\ \forall \omega' \in \mathcal{C}_{SG^{WP}_t - 1, \omega},\ \omega \neq \omega',
\end{equation}
and (\verb|NAC_rD_FD|)
\begin{equation} \label{EQ:NAC_rDFD}
	r^{D,FD}_{t, \omega} = r^{D,FD}_{t, \omega'} \quad \forall t \in \mathcal{T},\ \forall \omega \in \Omega,\ \forall \omega' \in \mathcal{C}_{SG^{WP}_t - 1, \omega},\ \omega \neq \omega',
\end{equation}

For BESS operational variables, we have (\verb|NAC_cV|)
\begin{equation} \label{EQ:NAC_cV}
	c_{t, \omega} = c_{t, \omega'} \quad \forall t \in \mathcal{T}, \ \forall \omega \in \Omega, \ \forall \omega' \in \mathcal{C}_{SG^{WP}_t - 1, \omega}, \ \omega \neq \omega',
\end{equation}
and (\verb|NAC_dV|)
\begin{equation} \label{EQ:NAC_dV}
	d_{t, \omega} = d_{t, \omega'} \quad \forall t \in \mathcal{T}, \ \forall \omega \in \Omega, \ \forall \omega' \in \mathcal{C}_{SG^{WP}_t - 1, \omega}, \ \omega \neq \omega',
\end{equation}
and (\verb|NAC_idV|)
\begin{equation} \label{EQ:NAC_idV}
	id_{t, \omega} = id_{t, \omega'} \quad \forall t \in \mathcal{T}, \ \forall \omega \in \Omega, \ \forall \omega' \in \mathcal{C}_{SG^{WP}_t - 1, \omega}, \ \omega \neq \omega',
\end{equation}
and (\verb|NAC_socV|)
\begin{equation} \label{EQ:NAC_socV}
	\soc_{t, \omega} = \soc_{t, \omega'} \quad \forall t \in \mathcal{T}, \ \forall \omega \in \Omega, \ \forall \omega' \in \mathcal{C}_{SG^{WP}_t - 1, \omega}, \ \omega \neq \omega'.
\end{equation}
and (\verb|NAC_rU_B|)
\begin{equation} \label{EQ:NAC_rUB}
	r^{U,B}_{t, \omega} = r^{U,B}_{t, \omega'} \quad \forall t \in \mathcal{T},\ \forall \omega \in \Omega,\ \forall \omega' \in \mathcal{C}_{SG^{WP}_t - 1, \omega},\ \omega \neq \omega',
\end{equation}
and (\verb|NAC_rD_B|)
\begin{equation} \label{EQ:NAC_rDB}
	r^{D,B}_{t, \omega} = r^{D,B}_{t, \omega'} \quad \forall t \in \mathcal{T},\ \forall \omega \in \Omega,\ \forall \omega' \in \mathcal{C}_{SG^{WP}_t - 1, \omega},\ \omega \neq \omega',
\end{equation}
\begin{note}
Observe that each of the nonanticipativity constraints \eqref{EQ:NAC_mUC}-\eqref{EQ:NAC_socV} might define the same constraint multiple times. This is not necessary in the implementation of the model, and in fact, we do not include them multiple times in our implementation.
	
Moreover, it is possible to reduce the number of constraints further. For example, let $a_{\omega}$ be a decision variable of stage $sg$ and consider the cluster $\mathcal{C}_{sg, \omega_1} = \{ \omega_1, \omega_2, \omega_3 \}$. We just need the following nonanticipativity constraints 
$$
a_{\omega_1} = a_{\omega_2}, \ a_{\omega_2} = a_{\omega_3},
$$
instead of the full set of nonanticipativity constraints (without repetitions)
$$
a_{\omega_1} = a_{\omega_2}, \ a_{\omega_1} = a_{\omega_3}, a_{\omega_2} = a_{\omega_3}.
$$

In practise, a presolver might also take care of this if not implemented efficiently.	
\end{note}


\section{Objective Function}
\label{SEC:ObjectiveFunction}

We define the objective function Expected Energy Community Social Welfare (EECSW) as follows:
\begin{subequations}
	\begin{align}
    EECSW & := \sum_{t \in \mathcal{T}} \sum_{\omega \in \Omega} \Pi_{\omega} \lambda^D_{t, \omega} (e^{DA+}_{t, \omega} - e^{DA-}_{t, \omega}) \label{EQ:DAincome} \\
&\;+\;
\sum_{t \in \mathcal{T}}
\sum_{\omega \in \Omega}
    \Pi_{\omega}\,
    \Bigl[
      \lambda^{R}_{t,\omega}
        \bigl(r^{U}_{t,\omega} + r^{D}_{t,\omega}\bigr)
      \;-\;
      \lambda^{R,\text{undel}}_{t,\omega}
        \bigl(r^{U,\text{undel}}_{t,\omega} + r^{D,\text{undel}}_{t,\omega}\bigr)
    \Bigr] \label{EQ:RESincome} \\
		& + \sum_{t \in \mathcal{T}} \sum_{\omega \in \Omega} \Pi_{\omega} \sum_{i \in \mathcal{I}_t} \lambda^{I}_{i, t, \omega} e^{IM}_{i, t, \omega} \label{EQ:IMincome} \\
		& + \sum_{t \in \mathcal{T}} \sum_{\omega \in \Omega} \Pi_{\omega} \lambda^{IB+}_{t, \omega} p^{IB+}_{t, \omega} \label{EQ:PosIBcollectRights} \\
		& - \sum_{t \in \mathcal{T}} \sum_{\omega \in \Omega} \Pi_{\omega} \lambda^{IB-}_{t, \omega} p^{IB-}_{t, \omega} \label{EQ:NegIBPayObligations} \\
		& - \sum_{t \in \mathcal{T}} \sum_{\omega \in \Omega} \Pi_{\omega} C^{FD} (afd^+_{t, \omega} + afd^-_{t, \omega}). \label{EQ:DemFlexPenalty}
	\end{align}
	\label{EQ:EECSW}
\end{subequations}

The second term \eqref{EQ:RESincome} represents the EC's expected income from the reserve market. The third term \eqref{EQ:IMincome} represents the EC's expected income from intraday markets. The forth term \eqref{EQ:PosIBcollectRights} represents the EC's expected positive imbalances collection rights. The fifth term \eqref{EQ:NegIBPayObligations} represents the EC's expected imbalance payment obligations. The last term \eqref{EQ:DemFlexPenalty} is a penalty on the use of demand flexibility.

\section{Complete Model Formulation}
\label{SEC:CompleteModelFormulation}

In this section, we summarise the complete model formulation. The objective is to maximize the social welfare of the energy community.

\begin{align*}
\notag & \text{max} & & EECSW \\
\notag & \text{s.t.} & & \eqref{EQ:FlexDemandPos}, \eqref{EQ:DailyDem}, \eqref{EQ:InterDem}, \eqref{EQ:FlexDemDisplace},  & \text{Flex. Demand Op.} \\
\notag & & & \eqref{dVdCPos}, \eqref{EQ:VPP_state_dV}, \eqref{EQ:VPP_state_cV}, \eqref{EQ:SOCV}, \eqref{EQ:socVOperLimits}, \eqref{EQ:SOCV_ini}, \eqref{EQ:SOCV_fin}, & \text{BESS Op.} \\
\notag & & & \eqref{EQ:RVarDef}, \eqref{EQ:RDcomp}  & \text{EC RM} \\
\notag & & & \eqref{EQ:DM_IM_1}, \eqref{EQ:IM_bounds_3-4}, & \text{EC IM} \\
\notag & & & \eqref{EQ:IBVarDef}, \eqref{EQ:IBDef}, \eqref{EQ:VPP_IB_MAX}, \eqref{EQ:VPP_IB_MIN}, & \text{EC IB} \\
\notag & & & \eqref{EQ:NAC_eIM}, \eqref{EQ:NAC_eIM2} & \text{NAC IM} \\
\notag & & & \eqref{EQ:NAC_pPIB}, \eqref{EQ:NAC_pNIB}, & \text{NAC Imbalances} \\
\notag & & & \eqref{EQ:NAC_var_fd}, \eqref{EQ:NAC_var_afd_p}, \eqref{EQ:NAC_var_afd_m}, & \text{NAC Flex. Demand} \\
\notag & & & \eqref{EQ:NAC_cV}, \eqref{EQ:NAC_dV}, \eqref{EQ:NAC_idV}, \eqref{EQ:NAC_socV}. & \text{NAC BESS Op.} \\
\end{align*}

\section{Summary of Stages}

In Table \ref{TABLE:stages-rv-dv}, we present a summary of the stages considered. Each stage can contain decision variables, observations and recourse variables. Note that, in terms of the optimization model, recourse variables of one stage are equivalent to decision variables of the next stage.

\begin{table}
\begin{tabular}{c c c c}
	Stage & Decision Variables & Observations & Recourse Variables \\
	4 & $e^{IM}_{2, t, \omega}$ & $\lambda^I_2 \in \R^{24}$ &  \\
	5 & $fd_1, afd^{\pm}_1, c_{1, \omega}, d_{1, \omega}, id_{1, \omega}, \soc_{1, \omega}, r^{U, B}_{1, \omega}, r^{U, FD}_{1, \omega}, r^{D, B}_{1, \omega}, r^{D, FD}_{1, \omega}$ & $W_1, PV_1, \lambda^{IB\pm}_{1, \omega} \in \R$ & $p^{IB+}_{1, \omega}, p^{IB-}_{1, \omega}$  \\
	6 & $fd_2, afd^{\pm}_2, c_{2, \omega}, d_{2, \omega}, id_{2, \omega}, \soc_{2, \omega}, r^{U, B}_{2, \omega}, r^{U, FD}_{2, \omega}, r^{D, B}_{2, \omega}, r^{D, FD}_{2, \omega}$ & $W_2, PV_2, \lambda^{IB\pm}_{2, \omega} \in \R$ & $p^{IB+}_{2, \omega}, p^{IB-}_{2, \omega}$ \\
	7 & $fd_3, afd^{\pm}_3, c_{3, \omega}, d_{3, \omega}, id_{3, \omega}, \soc_{3, \omega}, r^{U, B}_{3, \omega}, r^{U, FD}_{3, \omega}, r^{D, B}_{3, \omega}, r^{D, FD}_{3, \omega}$ & $W_3, PV_3, \lambda^{IB\pm}_{3, \omega} \in \R$ & $p^{IB+}_{3, \omega}, p^{IB-}_{3, \omega}$ \\
	8 & $fd_4, afd^{\pm}_4, c_{4, \omega}, d_{4, \omega}, id_{4, \omega}, \soc_{4, \omega}, r^{U, B}_{4, \omega}, r^{U, FD}_{4, \omega}, r^{D, B}_{4, \omega}, r^{D, FD}_{4, \omega}$ & $W_4, PV_4, \lambda^{IB\pm}_{4, \omega} \in \R$ & $p^{IB+}_{4, \omega}, p^{IB-}_{4, \omega}$ \\
	9 & $fd_5, afd^{\pm}_5, c_{5, \omega}, d_{5, \omega}, id_{5, \omega}, \soc_{5, \omega}, r^{U, B}_{5, \omega}, r^{U, FD}_{5, \omega}, r^{D, B}_{5, \omega}, r^{D, FD}_{5, \omega}$ & $W_5, PV_5, \lambda^{IB\pm}_{5, \omega} \in \R$ & $p^{IB+}_{5, \omega}, p^{IB-}_{5, \omega}$\\
	10 & $fd_6, afd^{\pm}_6, c_{6, \omega}, d_{6, \omega}, id_{6, \omega}, \soc_{6, \omega}, r^{U, B}_{6, \omega}, r^{U, FD}_{6, \omega}, r^{D, B}_{6, \omega}, r^{D, FD}_{6, \omega}$ & $W_6, PV_6, \lambda^{IB\pm}_{6, \omega} \in \R$ & $p^{IB+}_{6, \omega}, p^{IB-}_{6, \omega}$ \\
	11 & $fd_7, afd^{\pm}_7, c_{7, \omega}, d_{7, \omega}, id_{7, \omega}, \soc_{7, \omega}, r^{U, B}_{7, \omega}, r^{U, FD}_{7, \omega}, r^{D, B}_{7, \omega}, r^{D, FD}_{7, \omega}$ & $W_7, PV_7, \lambda^{IB+}_{7, \omega} \in \R$ & $p^{IB+}_{7, \omega}, p^{IB-}_{7, \omega}$ \\
	12 & $fd_8, afd^{\pm}_8, c_{8, \omega}, d_{8, \omega}, id_{8, \omega}, \soc_{8, \omega}, r^{U, B}_{8, \omega}, r^{U, FD}_{8, \omega}, r^{D, B}_{8, \omega}, r^{D, FD}_{8, \omega}$ & $W_8, PV_8, \lambda^{IB+}_{8, \omega} \in \R$ & $p^{IB+}_{8, \omega}, p^{IB-}_{8, \omega}$ \\
	13 & $fd_9, afd^{\pm}_9, c_{9, \omega}, d_{9, \omega}, id_{9, \omega}, \soc_{9, \omega}, r^{U, B}_{9, \omega}, r^{U, FD}_{9, \omega}, r^{D, B}_{9, \omega}, r^{D, FD}_{9, \omega}$ & $W_9, PV_9, \lambda^{IB\pm}_{9, \omega} \in \R$ & $p^{IB+}_{9, \omega}, p^{IB-}_{9, \omega}$\\
	14 & $fd_{10}, afd^{\pm}_{10}, c_{10, \omega}, d_{10, \omega}, id_{10, \omega}, \soc_{10, \omega}, r^{U, B}_{10, \omega}, r^{U, FD}_{10, \omega}, r^{D, B}_{10, \omega}, r^{D, FD}_{10, \omega}$ & $W_{10}, PV_{10}, \lambda^{IB\pm}_{10, \omega} \in \R$ & $p^{IB+}_{10, \omega}, p^{IB-}_{10, \omega}$ \\
    15 & $e^{IM}_{3, t, \omega}$ & $\lambda^I_3 \in \R^{12}$ &  \\
	16 & $fd_{11}, afd^{\pm}_{11}, c_{11, \omega}, d_{11, \omega}, id_{11, \omega}, \soc_{11, \omega}, r^{U, B}_{11, \omega}, r^{U, FD}_{11, \omega}, r^{D, B}_{11, \omega}, r^{D, FD}_{11, \omega}$ & $W_{11}, PV_{11}, \lambda^{IB\pm}_{11, \omega} \in \R$ & $p^{IB+}_{11, \omega}, p^{IB-}_{11, \omega}$ \\
	17 & $fd_{12}, afd^{\pm}_{12}, c_{12, \omega}, d_{12, \omega}, id_{12, \omega}, \soc_{12, \omega}, r^{U, B}_{12, \omega}, r^{U, FD}_{12, \omega}, r^{D, B}_{12, \omega}, r^{D, FD}_{12, \omega}$ & $W_{12}, PV_{12}, \lambda^{IB\pm}_{12, \omega} \in \R$ & $p^{IB+}_{12, \omega}, p^{IB-}_{12, \omega}$ \\
	18 & $fd_{13}, afd^{\pm}_{13}, c_{13, \omega}, d_{13, \omega}, id_{13, \omega}, \soc_{13, \omega}, r^{U, B}_{13, \omega}, r^{U, FD}_{13, \omega}, r^{D, B}_{13, \omega}, r^{D, FD}_{13, \omega}$ & $W_{13}, PV_{13}, \lambda^{IB\pm}_{13, \omega} \in \R$ & $p^{IB+}_{13, \omega}, p^{IB-}_{13, \omega}$ \\
	19 & $fd_{14}, afd^{\pm}_{14}, c_{14, \omega}, d_{14, \omega}, id_{14, \omega}, \soc_{14, \omega}, r^{U, B}_{14, \omega}, r^{U, FD}_{14, \omega}, r^{D, B}_{14, \omega}, r^{D, FD}_{14, \omega}$ & $W_{14}, PV_{14}, \lambda^{IB\pm}_{14, \omega} \in \R$ & $p^{IB+}_{14, \omega}, p^{IB-}_{14, \omega}$ \\
	20 & $fd_{15}, afd^{\pm}_{15}, c_{15, \omega}, d_{15, \omega}, id_{15, \omega}, \soc_{15, \omega}, r^{U, B}_{15, \omega}, r^{U, FD}_{15, \omega}, r^{D, B}_{15, \omega}, r^{D, FD}_{15, \omega}$ & $W_{15}, PV_{15}, \lambda^{IB\pm}_{15, \omega} \in \R$ & $p^{IB+}_{15, \omega}, p^{IB-}_{15, \omega}$ \\
	21 & $fd_{16}, afd^{\pm}_{16}, c_{16, \omega}, d_{16, \omega}, id_{16, \omega}, \soc_{16, \omega}, r^{U, B}_{16, \omega}, r^{U, FD}_{16, \omega}, r^{D, B}_{16, \omega}, r^{D, FD}_{16, \omega}$ & $W_{16}, PV_{16}, \lambda^{IB\pm}_{16, \omega} \in \R$ & $p^{IB+}_{16, \omega}, p^{IB-}_{16, \omega}$ \\
	22 & $fd_{17}, afd^{\pm}_{17}, c_{17, \omega}, d_{17, \omega}, id_{17, \omega}, \soc_{17, \omega}, r^{U, B}_{17, \omega}, r^{U, FD}_{17, \omega}, r^{D, B}_{17, \omega}, r^{D, FD}_{17, \omega}$ & $W_{17}, PV_{17}, \lambda^{IB\pm}_{17, \omega} \in \R$ & $p^{IB+}_{17, \omega}, p^{IB-}_{17, \omega}$ \\
	23 & $fd_{18}, afd^{\pm}_{18}, c_{18, \omega}, d_{18, \omega}, id_{18, \omega}, \soc_{18, \omega}, r^{U, B}_{18, \omega}, r^{U, FD}_{18, \omega}, r^{D, B}_{18, \omega}, r^{D, FD}_{18, \omega}$ & $W_{18}, PV_{18}, \lambda^{IB\pm}_{18, \omega} \in \R$ &  $p^{IB+}_{18, \omega}, p^{IB-}_{18, \omega}$ \\
	24 & $fd_{19}, afd^{\pm}_{19}, c_{19, \omega}, d_{19, \omega}, id_{19, \omega}, \soc_{19, \omega}, r^{U, B}_{19, \omega}, r^{U, FD}_{19, \omega}, r^{D, B}_{19, \omega}, r^{D, FD}_{19, \omega}$ & $W_{19}, PV_{19}, \lambda^{IB\pm}_{19, \omega} \in \R$ & $p^{IB+}_{19, \omega}, p^{IB-}_{19, \omega}$ \\
	25 & $fd_{20}, afd^{\pm}_{20}, c_{20, \omega}, d_{20, \omega}, id_{20, \omega}, \soc_{20, \omega}, r^{U, B}_{20, \omega}, r^{U, FD}_{20, \omega}, r^{D, B}_{20, \omega}, r^{D, FD}_{20, \omega}$  & $W_{20}, PV_{20}, \lambda^{IB\pm}_{20, \omega} \in \R$ & $p^{IB+}_{20, \omega}, p^{IB-}_{20, \omega}$ \\
	26 & $fd_{21}, afd^{\pm}_{21}, c_{21, \omega}, d_{21, \omega}, id_{21, \omega}, \soc_{21, \omega}, r^{U, B}_{21, \omega}, r^{U, FD}_{21, \omega}, r^{D, B}_{21, \omega}, r^{D, FD}_{21, \omega}$ & $W_{21}, PV_{21}, \lambda^{IB\pm}_{21, \omega} \in \R$ & $p^{IB+}_{21, \omega}, p^{IB-}_{21, \omega}$ \\
	27 & $fd_{22}, afd^{\pm}_{22}, c_{22, \omega}, d_{22, \omega}, id_{22, \omega}, \soc_{22, \omega}, r^{U, B}_{22, \omega}, r^{U, FD}_{22, \omega}, r^{D, B}_{22, \omega}, r^{D, FD}_{22, \omega}$ & $W_{22}, PV_{22}, \lambda^{IB\pm}_{22, \omega} \in \R$ & $p^{IB+}_{22, \omega}, p^{IB-}_{22, \omega}$ \\
	28 & $fd_{23}, afd^{\pm}_{23}, c_{23, \omega}, d_{23, \omega}, id_{23, \omega}, \soc_{23, \omega}, r^{U, B}_{23, \omega}, r^{U, FD}_{23, \omega}, r^{D, B}_{23, \omega}, r^{D, FD}_{23, \omega}$ & $W_{23}, PV_{23}, \lambda^{IB\pm}_{23, \omega} \in \R$ & $p^{IB+}_{23, \omega}, p^{IB-}_{23, \omega}$ \\
	29 & $fd_{24}, afd^{\pm}_{24}, c_{24, \omega}, d_{24, \omega}, id_{24, \omega}, \soc_{24, \omega}, r^{U, B}_{24, \omega}, r^{U, FD}_{24, \omega}, r^{D, B}_{24, \omega}, r^{D, FD}_{24, \omega}$ & $W_{24}, PV_{24}, \lambda^{IB\pm}_{24, \omega} \in \R$ & $p^{IB+}_{24, \omega}, p^{IB-}_{24, \omega}$ \\

\end{tabular}


\caption{Summary of stages, decision variables, random variable observations and recourse variables.}
\label{TABLE:stages-rv-dv}
\end{table}



\printbibliography



















\end{document}
